%% ma: L11-B14-0002
%% lop: 11
%% chuong: 4
%% bai: 14
%% dang: 11.14.1
%% loai: TL
%% muc: TH
%% dap_an: "-"
%% nguon: "tai-lieu-hoc-tap-mon-toan-lop-11-theo-chuong-trinh-giao-duc-pho-thong-2018.pdf (Lớp Toán Thầy Duy, Phạm Lê Duy, THPT Chu Văn An), Chương 4 Bài 5, Dạng 2, Ví dụ 2 (tr.523)"
%% kien_thuc: "Lớp 11 Bài 14 (phép chiếu song song)"
%% trang_thai: chinh-thuc
%% ghi_chu: "Đã duyệt 03/11/2026."
\begin{ex}
Cho hình chóp $S.ABC$. Gọi $I$ là trung điểm của $AB$.

a) Gọi $M,N,E$ lần lượt là trung điểm của $SA,SB,SC$. Tìm ảnh của $\triangle MNE$ trong phép chiếu song song trên mặt phẳng $(ABC)$ theo phương $SC$.

b) Tìm ảnh của tam giác $MNE$ trong phép chiếu song song trên $(ABC)$ theo phương trung tuyến $SI$ của $\triangle SAB$.
\loigiai{\sloppy a) Phương chiếu $SC$: $E\in SC$ nên hình chiếu của $E$ là $C$. Trong mặt phẳng $(SAC)$, đường thẳng qua $M$ song song $SC$ cắt $AC$ tại trung điểm $P$ của $AC$ (đường trung bình của $\triangle SAC$); tương tự đường thẳng qua $N$ song song $SC$ cắt $BC$ tại trung điểm $Q$ của $BC$. Ảnh của $\triangle MNE$ là $\triangle PQC$.

b) Phương chiếu $SI$: trong $(SAI)$, đường thẳng qua $M$ song song $SI$ cắt $AI$ tại trung điểm $M_1$ của $AI$; trong $(SBI)$, đường thẳng qua $N$ song song $SI$ cắt $BI$ tại trung điểm $N_1$ của $BI$; trong $(SCI)$, đường thẳng qua $E$ song song $SI$ cắt $CI$ tại trung điểm $E_1$ của $CI$. Ảnh của $\triangle MNE$ là $\triangle M_1N_1E_1$.}
\end{ex}
